\documentclass[11pt,a4paper]{article}

\usepackage[T1]{fontenc}
\usepackage{amsmath}
\usepackage{amssymb}
\usepackage{graphicx}
\usepackage{booktabs}
\usepackage[margin=25mm]{geometry}
\usepackage[authoryear,round]{natbib}
\usepackage{hyperref}

\hypersetup{
  colorlinks=true,
  linkcolor=black,
  citecolor=black,
  urlcolor=blue
}
\graphicspath{{figures/}}

\title{Validation of Newtonian long-bubble simulations in circular tubes}
\author{Vatsal Sanjay}
\date{September 2026}

\begin{document}

\maketitle

\begin{abstract}
A long gas bubble translating through a wetting liquid leaves an annular film
between its interface and the tube wall. We assess axisymmetric two-phase
simulations of this process against the Bretherton asymptote and a
finite-capillary-number representation of Taylor's measurements. Four
calculations at one maximum refinement level reproduce the reference flat-film
thickness within 2\% over $0.0054\lesssim Ca_b\lesssim0.068$. The agreement is
preliminary because no spatial refinement sequence has yet been completed.
We distinguish the flat film from the smaller local gap beneath the rear
meniscus and define a grid-sensitivity study using the four existing physical
inputs. This study tests whether the measured film thickness and bubble speed
are independent of spatial resolution before extending the capillary-number
range.
\end{abstract}

\section{Physical problem and reference limits}

Viscous stress draws liquid into the film around a translating bubble, while
surface tension connects the nearly cylindrical film to the curved front and
rear menisci. Let $R_t$ denote the tube radius, $R_v$ the volume-equivalent
bubble radius, $U_b$ the bubble speed, and $\overline U$ the imposed mean
liquid speed. We define
\begin{equation}
  H=\frac{b_{\mathrm{flat}}}{R_t}, \qquad
  Ca_b=\frac{\mu_c U_b}{\sigma}, \qquad
  Ca_{\mathrm{in}}=\frac{\mu_c\overline U}{\sigma},
  \label{eq:definitions}
\end{equation}
where $\mu_c$ is the continuous-phase viscosity and $\sigma$ is the surface
tension. The simulations prescribe $Ca_{\mathrm{in}}$, whereas the motion of
the front tip determines $Ca_b$. The reference laws must therefore be
evaluated at the measured bubble speed.

For a long inviscid bubble at small capillary and Reynolds numbers,
\citet{bretherton1961motion} matched the viscous dynamic meniscus to the
capillary outer shape and obtained
\begin{equation}
  H_B\simeq0.643(3Ca_b)^{2/3}\simeq1.337Ca_b^{2/3}.
  \label{eq:bretherton-film}
\end{equation}
The result also assumes negligible gravity, a perfectly wetting liquid, and
separation between the two dynamic menisci. The rounded coefficient 1.34 is
conventional, but the unrounded value matters when errors below one per cent
are discussed. A limiting-case test must show that
$H/Ca_b^{2/3}\rightarrow1.337$ as $Ca_b\rightarrow0$; a line with slope $2/3$
over a moderate range does not establish this limit.

Bretherton's second observable is the relative speed excess,
\begin{equation}
  W=\frac{U_b-\overline U}{U_b}
    \simeq1.29(3Ca_b)^{2/3}.
  \label{eq:speed-excess}
\end{equation}
If the uniform film is stationary in the laboratory frame, volume
conservation gives
\begin{equation}
  \frac{\overline U}{U_b}=(1-H)^2,
  \qquad W=1-(1-H)^2.
  \label{eq:volume-speed}
\end{equation}
The liquid flux in the simulated film must be measured before using
Eq.~\eqref{eq:volume-speed} as an exact identity. Finite gas viscosity,
meniscus interaction, and flow within the film can all modify it.

Pressure comparisons require a stated sign convention. Taking the positive
excess pressure from the liquid behind the bubble to the liquid ahead of it,
the summed front-and-rear dynamic contribution is
\begin{equation}
  \Delta p_{\mathrm{bubble}}
  \simeq4.52\frac{\sigma}{R_t}(3Ca_b)^{2/3}.
  \label{eq:pressure-drop}
\end{equation}
The coefficient 3.58 quoted in the abstract of
\citet{bretherton1961motion} is the leading-meniscus dynamic correction. It
should not replace 4.52 in a comparison across a complete long bubble.
Numerical pressures must be sampled outside both dynamic menisci using the
same sign and radius convention.

At finite capillary number, \citet{aussillous2000quick} represented the
circular-tube measurements of \citet{taylor1961deposition} by
\begin{equation}
  H_T=\frac{1.34Ca_b^{2/3}}{1+3.35Ca_b^{2/3}}.
  \label{eq:aq-correlation}
\end{equation}
This rational expression is an empirical fit introduced by Aussillous and
Qu\'er\'e, rather than an equation proposed by Taylor. It approaches
Eq.~\eqref{eq:bretherton-film} at small $Ca_b$ and accounts empirically for
geometrical confinement as the film thickens. Taylor measured the deposited
liquid fraction $m$ behind an advancing gas finger; for a uniform annular film,
$m=1-(1-H)^2$. Aussillous and Qu\'er\'e's high-speed experiment instead moved
a finite liquid column with air. Its inertial thresholds $Ca^*$ and $Ca^{**}$
are consequently not universal transition values for the finite gas bubbles
considered here.

\section{Numerical method and provenance}

The calculation solves the axisymmetric, incompressible two-phase
Navier--Stokes equations. A volume-of-fluid field tracks the gas--liquid
interface, surface tension supplies the capillary stress, and an embedded
boundary represents the circular tube wall. The no-slip wall is fully wetted
by the continuous phase. A Poiseuille profile drives the liquid at the inlet,
and the bubble translates without touching the wall.

Surface tension, $\mu_c$, and $R_v$ define the repeating scales. The present
calculations use
\begin{equation}
  \frac{R_t}{R_v}=0.7, \qquad
  \frac{\mu_g}{\mu_c}=0.01, \qquad
  \frac{\rho_g}{\rho_c}=0.001,
  \label{eq:properties}
\end{equation}
and the volume-radius Laplace number
\begin{equation}
  La_v=\frac{\rho_c\sigma R_v}{\mu_c^2}=1.
\end{equation}
The tube-based quantities are therefore
\begin{equation}
  La_t=0.7La_v, \qquad Re_t=La_tCa_b, \qquad We_t=La_tCa_b^2.
  \label{eq:dimensionless}
\end{equation}
At $Ca_b=1$, both $Re_t$ and $We_t$ are 0.7. Inertia must therefore be
measured during the extension to order-one capillary number, even though both
phases remain Newtonian.

The solver and numerical checks use the project-pinned Basilisk release
\texttt{v2026-08-30}. A reproducible calculation records this version, the
solver source, dimensionless inputs, compiler options, and analysis
definitions. The current results use a domain of length $16R_v$ and maximum
refinement level 10. The smallest available cell is then
$\Delta_{\min}=16R_v/2^{10}=0.015625R_v$.

\section{Preliminary flat-film comparison}

Figure~\ref{fig:film} compares the four completed calculations with
Eqs.~\eqref{eq:bretherton-film} and \eqref{eq:aq-correlation}. The plotted
data are also supplied in \nolinkurl{data/baseline-film.csv}.

\begin{figure}[t]
  \centering
  \includegraphics[width=\textwidth]{film-thickness-vs-ca.pdf}
  \caption{Flat-film thickness against the measured bubble capillary number.
  Symbols show the four simulations, the solid curve is the
  Aussillous--Qu\'er\'e representation of Taylor's data, and the dashed curve
  is the small-capillary-number Bretherton asymptote. The second panel shows
  the ratio to the finite-capillary-number correlation. The reference curves
  extend beyond the range of completed simulations.}
  \label{fig:film}
\end{figure}

\begin{table}[t]
  \centering
  \caption{Measured flat-film thickness and finite-capillary-number reference
  values. All four simulations use maximum refinement level 10.}
  \label{tab:baseline}
  \begin{tabular}{rrrrr}
    \toprule
    $Ca_{\mathrm{in}}$ & $Ca_b$ & $H$ & $H_T$ & $100(H/H_T-1)$ \\
    \midrule
    0.005 & 0.00538868 & 0.0366354 & 0.0373420 & $-1.892\%$ \\
    0.010 & 0.0112414  & 0.0567910 & 0.0575664 & $-1.347\%$ \\
    0.020 & 0.0239374  & 0.0861486 & 0.0870712 & $-1.060\%$ \\
    0.050 & 0.0680939  & 0.143663  & 0.143363  & $+0.209\%$ \\
    \bottomrule
  \end{tabular}
\end{table}

The flat-film estimator removes 30\% of the bubble's axial extent at each end,
takes the median interface radius within the remaining central region, and
averages the result over the final quarter of the usable snapshots. A linear
fit to the front-tip position over the same temporal window gives $U_b$. This
procedure is designed to exclude both dynamic menisci. The existence of a
spatial plateau and sensitivity to the trimming and averaging windows have not
yet been quantified.

The global minimum gap near the rear meniscus is smaller. Its deviations from
$H_T$ are approximately $-32.3\%$, $-25.9\%$, $-23.7\%$, and $-18.4\%$ for
the four cases in Table~\ref{tab:baseline}. These values describe a different
observable: rear capillary waves can produce a local minimum below the
uniform film, as found in the rear-meniscus calculations of
\citet{giavedoni1999rear}. A steady minimum gap does not by itself establish a
steady flat film, pressure drop, bubble length, or complete travelling shape.

At level 10, the measured flat films span at most 1.64, 2.54, 3.86, and 6.44
smallest cells. These are geometric estimates, rather than verified counts of
the adaptive leaves within each film. No case presently has a completed
refinement sequence. The observed agreement within 2\% is therefore a
single-resolution comparison and cannot yet be interpreted as numerical
accuracy.

\section{Grid sensitivity and extension in capillary number}

The first grid study holds the four existing input values
$Ca_{\mathrm{in}}=0.005$, 0.01, 0.02, and 0.05 fixed. Each physical input is
examined at maximum levels 9, 10, 11, and 12, where level 10 is the existing
calculation. Selected low-capillary-number cases may require finer stages if
these four levels do not reach an asymptotic range. Table~\ref{tab:grid}
records the design; only its level-10 column contains results.

\begin{table}[t]
  \centering
  \caption{Grid-sensitivity matrix. ``Pending'' denotes a calculation under
  investigation and carries no implied result.}
  \label{tab:grid}
  \begin{tabular}{rccccc}
    \toprule
    $Ca_{\mathrm{in}}$ & Level 9 & Level 10 & Level 11 & Level 12 & Finer \\
    \midrule
    0.005 & pending & existing & pending & pending & conditional \\
    0.010 & pending & existing & pending & pending & conditional \\
    0.020 & pending & existing & pending & pending & conditional \\
    0.050 & pending & existing & pending & pending & conditional \\
    \bottomrule
  \end{tabular}
\end{table}

The same inlet capillary number and physical properties are retained across a
grid sequence. Both $H$ and the resulting $Ca_b$ must converge; retuning the
inlet speed would hide part of the numerical response. The correlation is
evaluated at the measured $Ca_b$ on every grid. Bubble volume, length, flat-film
uniformity, and temporal drift accompany the two primary observables. At
least three solutions in an asymptotic range are required before assigning a
convergence order. Time-step and pressure-solver tolerances, domain length,
boundary clearance, bubble length, and gas viscosity ratio are separate tests
and should not be folded into the spatial sequence.

Adaptive refinement currently responds to interface, curvature, velocity, and
strain-rate indicators. A feature-based criterion that guarantees a minimum
number of cells across the wall--interface gap may reduce the cost of finer
stages. This remains a proposed numerical method. Its feature definition,
refinement floor, spatial band, interaction with the embedded boundary, and
volume-conservation tests must be fixed before it can contribute evidence to
the grid study.

Once spatial sensitivity is known, the bubble calculations can extend from
the small-$Ca_b$ limit to $Ca_b=1$. Equation~\eqref{eq:aq-correlation} supplies
a convenient film-thickness comparison but does not replace independent
benchmarks. The axisymmetric Stokes solutions of
\citet{reinelt1985penetration} provide deposited fraction, finger radius, nose
shape, and front pressure response at finite $Ca_b$.
\citet{giavedoni1997axisymmetric} and \citet{giavedoni1999rear} give
complementary front and rear reference problems, including streamline
topology and rear interfacial undulations. Their plane-channel results must not
be mixed with the circular-tube geometry.

\section{Present conclusion}

The present calculations support one bounded statement. Four level-10
simulations of Newtonian gas bubbles reproduce the flat-film value of
Eq.~\eqref{eq:aq-correlation} within 2\% over
$0.0054\lesssim Ca_b\lesssim0.068$. The minimum gap beneath the rear meniscus
does not test the same quantity. Spatial convergence, the limiting Bretherton
coefficient, speed excess, meniscus profiles, and pressure response remain
unestablished. The pending refinement matrix is designed to quantify the first
of these uncertainties before the comparison is extended through $Ca_b=1$.

\bibliographystyle{plainnat}
\bibliography{references}

\end{document}
